\documentclass{article}
\usepackage{amsmath, amsthm, amsfonts, amssymb}
\usepackage{hyperref}
\theoremstyle{plain}
\newtheorem{thm}{Theorem}
\numberwithin{thm}{section}

\theoremstyle{plain}
\newtheorem{prop}{Proposition}
\numberwithin{prop}{section}

\theoremstyle{definition}
\newtheorem{defn}{Definition}
\numberwithin{defn}{section}

\theoremstyle{remark}
\newtheorem*{rem}{Remark}

\newtheorem*{cor}{Corollary}

%\numberwithin{equation}{section}
\newcommand{\op}{\prime}
\begin{document}
\title{On a Generalisation of the Definite Integral}
\author{Henri Lebesgue}
\maketitle

An English translation of Lebesgue's 
\href{https://fermatslibrary.com/s/on-a-generalization-of-the-definite-integral}
{original paper} was used for this note.

\begin{enumerate}
\item An integral has two roles in calculus. It tells the area under a curve
between two points and it is an anti-derivative. If $y = f(x)$ is a curve then
the area under it between points $x = a$ and $x = b$ is given by
\[
\int_a^b f(x)dx.
\]
If $f = dg/dx$ then the above expression can be evaluated to $g(b) - g(a)$.
If we denote the differentiation by $D$ and integration by $I$ then if $Df = g$
then $f = I(g)$ so that $I = D^{-1}$. That is, integration is inverse of
differentiation and that is why it is called as an anti-derivative. The
equation
\begin{equation}\label{e1}
\int_a^b \frac{dg}{dx}dx = g(b) - g(a)
\end{equation}
was given the name \emph{the fundamental theorem of calculus} to emphasise its
important role in relating the two most important operations in the subject.

\item It was known since Newton's time that the left hand side of equation 
\eqref{e1} could be evaluated for an integrand $f$ even if we did not know a
function $g$ such that $f = Dg$. If $f$ was continuous, one could divide the
area under the curve $y = f(x)$ in slender enough rectangles and sum their 
areas. The operation of finding areas in this manner was called 
\emph{quadrature}. It was possible to find the quadrature even if one cannot
find the anti-derivative. For example, one could find the area under $y = 
\exp(-x^2)$ between $x = 0$ and $x = 1$ without knowing its anti-derivative.

\item However, there are functions $f$ and $g$ related as $f = Dg$ for which
the quadrature between $x=a$ and $x=b$ is not equal to $g(b) - g(a)$. For
functions like these, the fundamental theorem of calculus seems to fail. Two 
famous examples of functions for which this happened were Volterra's function
and Devil's staircase. Both of them have an anti-derivative but a quadrature
cannot be computed using Riemann's algorithm. Lebesgue sought to find a remedy
to this situation.

\item Lebesgue started with Riemann's definition of an integral of an 
increasing continuous function $f$ in the interval $[a, b]$
\begin{equation}\label{e2}
\int_a^b f(x)dx = 
\lim_{N \rightarrow \infty}\sum_{n=1}^N f(\xi_i)(x_{i+1} - x_i)
\end{equation}
where $a = \xi_1 < \xi_2 < \ldots \xi_{N+1} = b$ is a division of the domain
of the function into disjoint sub-intervals. As $f$ is an increasing continuous
function, if $x \in [x_i, x_{i+1}]$ then $y = f(x) \in [m_i, m_{i+1}]$, where
$m_i = f(x_i)$. Lebesgue noticed that the converse statement, that is, $y \in
[m_i, m_{i+1}] \Rightarrow x \in [x_i, x_{i+1}]$ is also true.

\item Continuing with this observation, Lebesgue divided the $y$ axis between
$y = m = f(a)$ and $y = M = f(b)$ at points $m = m_1 < m_2 \ldots < m_P = M$.
Define the sets $E_i = \{x \in E_i \Leftarrow y \in [m_i, m_{i+1}]\}$. Let 
$\lambda_i$ be the measure of $E_i$. For our choice of $f$, $\lambda_i$ is the
length of the interval $[x_i, x_{i+1}]$. If
\begin{equation}\label{e3}
L(a, b) = \lim_{P \rightarrow \infty}\sum_{i=1}^P m_i\lambda_i
\end{equation}
then $L(a, b)$ is the integral of the function between $x=a$ and $x=b$. It is
the same as the one calculated by Riemann's procedure in equation \eqref{e2}.
Lebesgue called a function $f$ integrable if the limit in \eqref{e3} existed.

\item Measure theory was introduced by 1900. Lebesgue used the following ideas
from measure theory:
\begin{enumerate}
\item The interval $(a, b)$ can be written as a union of intervals in an
infinite number of ways. Let $D^j = \{I^j_1, I^j_2\, \ldots\}$ be a family of
intervals covering $(a, b)$, that is $(a, b) \subseteq \cup_{n \ge 1}I^j_n$ so 
that the measure of $(a, b)$ is $\mu((a, b)) \le \sum_{n\ge 1}\mu(I^j_n)$. He 
defined,
\begin{equation}\label{e4}
\mu((a, b)) = 
\inf_{j \ge 1}\left\{\sum_{n \ge 1}\mu(I^j_n)\;;\; I^j_n \in D^j\right\}.
\end{equation}
\item A subset $E$ of $(a, b)$ is measurable if 
\begin{equation}\label{e5}
\mu(E) + \mu(E^c) = \mu((a,b)).
\end{equation}
\item If $E_i$ are measurable then so is their union and intersection.
\item If $E_i$ are pairwise disjoint then the measure of their union is the
sum of their measures.
\end{enumerate}
The unions and intersections were assumed to be countable.

\item In terms of a measure, the integral of equation \eqref{e3} can be 
written as 
\begin{equation}\label{e6}
L(a, b) = \lim_{P \rightarrow \infty}\sum_{i=1}^Pm_i\mu(E_i),
\end{equation}
where $E_i = \{x \in (a, b)\;:\; m_i < f(x) \le m_{i+1}\}$. 

\item From equation \eqref{e6} it follows that if a function $f$ lies between
limits $A$ and $B$ and if we divide the range of $f$ as $A = m_1 < m_2 \ldots
m_{P+1} = B$ and if $E_i = \{x \in (a, b)\;:\; m_i < f(x) \le m_{i+1}\}$ are
all measurable then the integral of $f$ exists. Lebesgue called this an
\emph{integrable} function. Note that there are no restrictions on $f$ except
that its range is bounded.

\item Instead of vieweing $L(a, b)$ in equation \eqref{e6} as a limit of a sum,
we can express it as
\begin{equation}\label{e7}
m_1\mu(E_1) + \cdots + m_P\mu(E_P) \le L(a, b) \le m_2\mu(E_1) + \cdots +
m_{P+1}\mu(E_P).
\end{equation}
The integral $L(a, b)$ is thus bounded between lower and upper sums in the 
sense of Darboux. Lebesgue called a function for which \eqref{e7} is true as 
\emph{summable}. An integral of a summable function lies between the upper and
lower integrals.

\item Lebesgue then claimed that if an integrable function is summable in the
sense of Riemann then the integral is the same in both definitions.

Consider a function $y = f(x)$. If $R(a, b)$ is the Riemann integral with 
$R_l$ and $R_u$ as the lower and upper Riemann sums then 
\begin{equation}\label{e8}
R_l \le R(a, b) \le R_u
\end{equation}
implies that for a subdivision $m_1, \ldots, m_P$ of the range of $f$ the 
pre-image of $(m_i, m_{i+1})$ are measurable subsets $E_i$ of the interval 
$(a, b)$ such that 
\begin{equation}\label{e9}
L_l \le L(a, b) \le L_u,
\end{equation} 
where $L_l = m_1\mu(E_1) + \cdots + m_P\mu(E_P)$, $L_u = m_2\mu(E_1) + \cdots 
+ m_{P+1}\mu(E_P)$ are the lower and upper Lebesgue sums and $L(a, b)$ is the 
Lebesgue integral.

\item If a function is summable in Riemann sense, that is if $R(a, b)$ can be
squeezed between upper and lower Riemann sums then the function has 
discontinuities at points whose union is a set of measure zero. By omitting
these points from $(a, b)$, the function is continuous on what remains. Such
a function is summable in the Lebesgue sense. As an example, let $f$ and $\phi$
be functions such that $\phi$ is not zero everywhere in $(a, b)$ and $f$ be 
summable. Let $Z$ be a subset of $(a, b)$ of measure zero. Now define a 
function $g$ as
\[
g(x) = \begin{cases} f(x) + \phi(x) & \;\forall\; x \in Z \\
f(x) & \;\forall\; x \in (a, b) \setminus Z.
\end{cases}
\]
Lebesgue claims that $g$ is summable but not in the Riemann sense.

For a sub-interval $(x_i, x_{i+1})$ of $(a, b)$ howsoever small, the denseness
of $Z$ guarantees that the upper Riemann sum is $(f(\xi) + \phi(\xi))(x_{i+1}
- x_i)$ while the lower Riemann sum is $f(\xi)(x_{i+1} - x_i)$, where $\xi \in
(x_i, x_{i+1})$, and the two sums will always differ because $\phi$ is not 
everywhere zero.

An example of this arrangement is $f = 0$ over $(0, 1)$, 
\[
\phi(x) = \begin{cases} 1 & x \in \mathbb{Q} \\
 0 & x \in (0, 1) \setminus \mathbb{Q}
\end{cases}
\]
and $g = f + \phi$ is the Dirichlet function. $\int_0^1 g = 0$ in the Lebesgue
sense but the integral does not exist in the Riemann sense.

\item If $f$ and $\phi$ are summable functions, then so is $f + \phi$. 
Furthermore, the integral of $f + \phi$ is a sum of integrals of $f$ and $\phi$.
If we show that $x^n$ are summable functions then so are all polynomials.

\item If $\{f_n\}$ is a sequence of summable functions then so is their limit.
Thus, the limits of polynomials, which are all continuous functions are 
summable.

\item If a derivative is bounded, then it is summable. Therefore, 
\[
\psi(y) = \int_a^y D\phi(x)dx
\]
then $\psi(y)$ exists and is the anti-derivative of $\phi$.

\item Lastly, if the derivatives $Df, D\phi, D\psi$ are merely bounded without
being continuous, the arc length can be calculated as a Lebesgue integral. If
$\phi$ and $\psi$ are zero, then 
\[
\int_a^b |Df(t)|dt
\]
is the total variation of the function $f$ between $a$ and $b$. This too can
be calculated as a Riemann integral as long as it is bounded even though it
oscillates wildly.

\item Even in situations where $Df$ does not exist and we can only have the 
Dini derivatives, a Lebesgue integral can be calculated.

\end{enumerate}
\end{document}

